\section{Related Work}\label{sec:related_work}

\subsection{Sensor-Based Process Monitoring in CNC Machining}

Sensor-based monitoring of CNC machining processes has been an active research area for over three decades. Teti et al.~\cite{teti2010advanced} reviewed force, vibration, acoustic emission, and current sensors for machining operations. Byrne et al.~\cite{byrne1995tool} established early tool condition monitoring using multiple sensor modalities. More recently, Nath~\cite{nath2020integrated} reviewed integrated sensor systems for advanced manufacturing, and Kuntoğlu et al.~\cite{kuntoglu2021review} surveyed sensor-based approaches for Industry 4.0 CNC monitoring. Low-cost MEMS-based inertial measurement units (IMUs) have expanded the range of deployable sensor configurations~\cite{mohring2012imu}, enabling dense instrumentation of machine structures at modest cost.

A recurring finding is that sensor signals during machining reflect both process parameters (feed rate, depth of cut, tool engagement) and machine-specific factors (structural dynamics, bearing condition, drive characteristics). While most monitoring approaches treat the combined signal as the object of interest, our work decomposes these contributions through air-cut referencing.

\subsection{Sensor Fusion for Manufacturing}

Multi-sensor data fusion has demonstrated advantages over single-sensor approaches for machining process characterization~\cite{cai2017sensor, sick2002online}. Combining heterogeneous sensor types, such as accelerometers with current sensors, captures complementary aspects of the process. Feature-level fusion, in which statistical features from individual sensor channels are concatenated into a joint feature vector, is the most common strategy and is the approach adopted in this work. Zhou et al.~\cite{zhou2023multisensor} reviewed multi-sensor fusion architectures for intelligent manufacturing, and deep-learning-based fusion has been applied to larger datasets~\cite{zhao2019deeplearning}.

\subsection{Signal Processing and Feature Extraction}

Time-domain statistical features (mean, variance, RMS, kurtosis, skewness) and frequency-domain features (spectral energy, dominant frequency, power spectral density) are widely used for characterizing machining signals~\cite{dimla2000sensor, kuo2000spectral, albadour2011wavelet}. Rmili et al.~\cite{rmili2016automatic} classified tool wear automatically from statistical and spectral features. Our temporal binning approach is related to short-time feature extraction, in which signals are divided into temporal segments and statistics are computed per segment.

\subsection{Machine Learning for Toolpath and Process Classification}

Random Forests~\cite{breiman2001randomforests}, Support Vector Machines~\cite{widodo2007svm, sun2004svm}, and neural networks have all been applied to machining process classification. Bustillo et al.~\cite{bustillo2022smart} demonstrated smart manufacturing applications of ensemble methods. Wang et al.~\cite{wang2018deep} and Cao et al.~\cite{cao2024deep} reviewed deep learning approaches for manufacturing quality prediction. Bergs et al.~\cite{bergs2020digital} positioned machine learning within the digital twin framework for manufacturing.

Recent work on time-series classification offers alternatives to the bag-of-statistics paradigm used here. Dempster et al.~\cite{dempster2020rocket} introduced ROCKET (Random Convolutional Kernel Transform), which applies random convolutional filters directly to time series and reports competitive accuracy at low computational cost. Fawaz et al.~\cite{fawaz2020inceptiontime} proposed InceptionTime, a deep learning architecture for time-series classification. We compare our temporal binning approach against ROCKET-style features in Section~\ref{sec:results}.

\subsection{Transfer Learning and Domain Adaptation}

The cross-machine generalization problem, in which models trained on one machine fail to transfer to another, has been addressed through transfer learning~\cite{pan2010transfer, li2020transfer} and domain adaptation~\cite{jiao2020domain}. These approaches typically require some labeled data from the target domain. Air-cut referencing is complementary. By normalizing against a machine-specific baseline, it may produce features that are inherently more transferable, a hypothesis that requires cross-machine validation (Section~\ref{sec:discussion}).

\subsection{Air-Cut and No-Load Testing}

Air cutting (also called dry running or no-load testing) is standard practice in machine tool acceptance testing and calibration~\cite{iso230_2012}. It characterizes spindle dynamics, axis positioning accuracy, and thermal behavior without cutting forces. Altintas~\cite{altintas2012manufacturing} describes machine dynamics characterization through no-load testing. In tool-condition monitoring, subtracting a no-load (idle) spindle- or axis-motor current/power baseline to isolate the cutting-load component is itself a well-established practice~\cite{teti2010advanced, dimla2000sensor, byrne1995tool}. For this reason, we frame our near-constant spindle-current result (Section~\ref{sec:physics}) as a known limitation of current-based monitoring under shallow cuts rather than a novel finding. Air-cut referencing generalizes the no-load-baseline idea in two ways: from a single current channel to the full multi-modal, temporally-binned feature vector (110 channels $\times$ per-phase statistics), and from load estimation to the removal of machine- and \emph{workspace-position}-specific offsets prior to classification. To the best of our knowledge, no prior work has systematically used air-cut data as a \emph{reference baseline across the full sensor array} for isolating process-specific signal components to enable toolpath classification, with condition discrimination demonstrated but not yet validated (Section~\ref{sec:damage}).

\subsection{CNC Security and Manufacturing Integrity}

CNC manufacturing security has gained attention as machines become networked~\cite{monostori2016cyber}. Sturm et al.~\cite{sturm2017cyber} and Zeltmann et al.~\cite{zeltmann2016manufacturing} demonstrated that malicious modifications to G-code or process parameters can introduce defects undetectable by conventional inspection. Yampolskiy et al.~\cite{yampolskiy2018security} surveyed security threats across manufacturing modalities. Belikovetsky et al.~\cite{belikovetsky2017drowned} showed that sensor-based verification can detect sabotaged manufacturing processes. Our toolpath fingerprinting methodology supports this application by providing sensor-based toolpath verification.
